\section{Literature review}
\label{sec:lit_review}

The relationship between media and government has been extensively studied in the fields of economics and political science, and has been identified as a critical component of electoral accountability in Latin America \citep{Audits-Ferraz-QJE, PolInfo-Enriquez-JEEA, Marshall}. A well-functioning media wields significant power to hold government actions accountable and inform the public, thereby fostering transparency and responsiveness \citep{Stromberg-MediaCoverage, Account-Boix-JoLEO}. 

Crucially, incumbent governments are fully aware of the impact that media and the broader information environment have on their chances of remaining in power, prompting them to devote substantial resources to control the information reaching the public \citep{Autocracy-Guriev-JPE}. In the following sections, I will outline various ways in which governments can influence the media and manipulate the information environment, considering both authoritarian and democratic contexts, as well as the particular Mexican context.

\subsection{Authoritarian Regimes}

In authoritarian contexts, governments exercise control over the media in two main ways. The first is by allocating substantial resources to control the information environment through censorship. For example, in China, \cite{Army-Chen-Book} and \citep{CensorshipCollective-King-APSR} document that each internet content provider employs thousands of censors to ensure compliance with government regulations, while the government itself deploys an extensive force of over 300,000 cyber-police to oversee censorship efforts. As described by \cite{CensorshipChina-King-Science}, this represents the largest recorded suppression of free expression in history. Notably, their seminal work shows that this censorship primarily targets posts indicative of collective action, while government criticism is generally allowed.

The second method of control is through personnel management. For instance, \cite{Plus-Anne-PostComm} note that the Chinese government directly controls the appointment of managers and editors within the media, while \cite{PersonnelMedia-Lee-IntJPPoli} highlight that dictators selectively promote aligned journalists to shape the content of news reports. Finally, some theoretical and empirical evidence suggest that in resource-scarce dictatorships, governments may be more inclined to permit greater press freedom as a means of gathering information about their bureaucrats \citep{PoorDict-Guriev-APSR} and power-sharing elites \citep{ElitesMedia-Tung-PolComm}. This approach allows the government to monitor and control these groups, helping to prevent crises and maintain stability.

\subsection{Democracies}

In well-functioning democracies, evidence of such practices is more limited. However, while these tactics are more nuanced, governments and elites still employ them to gain control over the information environment and influence media reporting.\footnote{By \textit{well-functioning democracies}, I refer to countries with a comprehensive set of democratic mechanisms, such as a constitution, regular elections, a free press, and opposition parties, among others. Given that the level of democratization varies widely across countries, I provide examples and differentiate between developing nations and industrialized democracies.} In this context, one common strategy employed by leaders is to delegitimize the media. There is substantial evidence that this tactic affects how individuals perceive media bias and damages the media’s overall reputation \citep{PercMedia-Egelholfer-JoComm, PolNews-Smith-JoPressPol}, which could in consequence affect their reporting. Moreover, this approach is not limited to developing countries \citep{JournIndia-Bhat-Dig, ThugMex-VictorHugo-JourPract}; it is also widespread in industrialized democracies \citep{KoreaJour-Shin-Journalism, PressCriticism-Carlson-Journalism}. In recent years, with the rise of populism, the use of this strategy has become more pronounced, further shaping public perception of the media \citep{PolEconPop-Guriev-JEL}.

Another tactic employed by governments to control the media is bribery. \cite{Peru-Macmillan-JEP} document how Montesinos, the intelligence officer under Fujimori, paid millions of dollars in bribes to media owners, opposition politicians, and judges. In some cases, these bribes were used to gain full control over media broadcasts for specific periods. Notably, and in support of my argument about the significance of the information environment, the amounts paid to media owners were exponentially higher—100 times greater—than those paid to judges and politicians.

Politicians and governments can also shape media reporting through legal mechanisms that require minimal resources. These include laws that restrict free speech, limit access to sources, or regulate media licensing \citep{PolicyBrief-MediaRegulation}. An interesting example of this subtle control is the use of kisha clubs in Japan. These press clubs are informal, invitation-only gatherings between government officials and select members of the media. They provide exclusive access to sources, which has raised concerns about press independence \citep{Kisha-Freedom}. Critics argue that the uniformity in news coverage across different newspapers stems from their reliance on the same controlled sources provided through these clubs \citep{Kisha-UoNott}.  

Related with laws that restrict free speech, \cite{MexicanPressStanigAJPS} analyzes cross-state variations in the strictness of speech regulations in Mexico to show that media outlets report less on government corruption in regions where speech is more tightly controlled. This trend is largely driven by journalists' fear of prosecution or defamation lawsuits. Using a similar approach, \cite{Journalism-Salazar-RevMex} finds that this effect is mediated by the level of judiciary independence in each state.

\section{The Mexican Case}

To adequately understand the relationship between the media and the Mexican government—and the broader necessity for governments to control the information environment—it is necessary to analyze the history of Mexico through three distinct periods. The first period, often referred to as the \textit{perfect dictatorship,} encompasses the 70-year reign of the Institutional Revolutionary Party (PRI) until the 2000s. The second period marks a transition initiated by the victory of the National Action Party (PAN) in the 2000 presidential election under President Vicente Fox followed by the PRI's return to power in 2012. The final period is the shift to Andrés Manuel López Obrador's (AMLO) government in 2018.

In each of these periods, the evolving power dynamics between the state and the media and their respective behaviors in managing the information environment provide valuable insights. These historical shifts can help us theorize the significance of changes in government advertising spending and its impact on media operations. Additionally, by reviewing the last 18 years of the violence context in Mexico, I can further argue how these changes in advertising might influence media reporting on this critical and salient issue of violence, affecting both the content and the depth of the coverage.

\subsection{The perfect dictatorship}

In 1970, author Mario Vargas Llosa famously described the Mexican government as a \textit{perfect dictatorship}, arguing that it had struck agreements with elites and was financing opposition groups to silence critics \citep{DictaduraPerf-Llosa}. Initially, this claim might have seemed exaggerated, but history now shows how the PRI used authoritarian techniques, such as use of the army and police, to suppress any dissent or insurrection. They controlled the media and elites, manipulating public opinion to maintain regime stability. This description was particularly relevant coming two years after the Tlatelolco massacre in 1968, where the Mexican army killed hundreds of protesting students—an event known internationally. Less known, however, was the subsequent 'Dirty War' which followed these events. Recent government reports indicate that approximately 2,600 individuals were direct or indirect victims of government repression during this period. The methods employed ranged from torture to mass executions \citep{ComisionParaVerdad}.

For media owners, the situation was markedly different. As \cite{FourthState-Lawson-Book} describes, their relationship with the regime was one of collusion rather than coercion. In his extensive documentation of the interactions between the government and the media—referred to as the 'Fourth State'—he provides both anecdotal and historical evidence of media owners maintaining business-like relationships with the PRI. These relationships were underpinned by mutual agreements to uphold the status quo, allowing the government to manage media narratives without resorting to violent repression.

However, the stability of such regimes is not merely based on an abstract idea of 'status quo maintenance'. The PRI also strategically deployed resources to keep the media loyal. This included tactics like tax exemptions, broadcasting concessions, protection from competition, lucrative business opportunities outside of media and \textit{chayote}, which was a particular and common form of bribery to reporters. For instance, \cite{FourthState-Lawson-Book} documents that Televisa remained an outright monopoly of TV in Mexico from 1972 to the 1990s. The firm received over this period of time the first television licenses, control of government built satellites, among other things. Moreover, when Televisa lost its monopoly power to dominant privatized entrant TV Azteca, Carlos Salinas de Gortari made sure that his brother was a silent partner in the deal. 

For smaller or independent media outlets, PRI leaders strategically utilized advertising spending as a tool to secure favorable coverage. Although specific data from that period remains undisclosed, the extent of this practice is evident without needing substantial quantitative evidence. This is exemplified by a notable incident involving President López Portillo, who candidly admitted to the strategy. On Freedom of the Press Day, ironically, he famously declared, \textit{'Pay for them to beat up on me? Gentlemen, I think not!'} while announcing the withdrawal of advertising spending to independent media \citep{PrensaVendida-Rod}.  

For independent or media that dissented with the government, other strategies were employed, some of which I already mentioned in past sections, such as manipulation of access, harassment and outright repression \citep{FourthState-Lawson-Book}. Legal measures such as revoking licenses and blacklisting media outlets from receiving official government information were common. Additionally, undisclosed firings of reporters who expressed dissent were widespread. Finally, from 1980 to 1996, although not all tied to government repression as cartel violence was also starting, approximately 60 journalists were killed, with some particularly significant cases such as the murder of Manuel Buendía. His death in particular was linked to the government, as he was looking into drug-related corruption within the federal police \citep{FourthState-Lawson-Book}. 

This period illustrates that the Mexican government has historically used state resources to manipulate the information environment, underscoring the critical importance of such control in a country like Mexico. Moreover, on the period we are analyzing (2012 to 2018), PRI came back into power after twelve years of PAN governments, therefore it is prudent to look back at the past behavior of the party and underscore the importance that media represented for them. The government's actions clearly show a strategic awareness of the power of media influence, which I argue continued to the present.

\subsection{Democratic transition}

Starting in the 1990s, \cite{FourthState-Lawson-Book} documents a significant shift in media behavior that cannot be fully explained by an increase in high-profile, damaging events for the government or changes in Mexico's broader institutional context. Instead, this shift is attributed to the emergence of a wave of '\textit{independent}' news outlets, such as El Norte (now Reforma), El Economista, El Financiero, La Jornada, etc. These outlets played a pivotal role in expanding coverage of critical issues such as massacres, corruption scandals, assassination conspiracies, and, notably, drug-related violence.\footnote{Examples of these events are described in detail by \cite{FourthState-Lawson-Book}, here are the ones I consider more relevant for each topic: For massacres the one in \href{https://www.theadvocatesforhumanrights.org/Res/guerrero_1995_2.pdf}{Aguas Blancas} and the one of \href{https://www.amnesty.org/en/wp-content/uploads/2021/06/amr410431998en.pdf}{Acteal}, for corruption scandals such as the campaign funding of \href{https://elpais.com/diario/2000/10/25/internacional/972424820_850215.html}{Roberto Madrazo}, assassination conspiracies such as the one from PRI's presidential candidate \href{https://www.jornada.com.mx/noticia/2024/03/23/politica/al-cumplirse-30-anos-del-asesinato-de-colosio-el-caso-esta-abierto-ahora-se-habla-de-una-accion-concertada-4892}{Luis Donaldo Colosio} and PRI's party leader \href{https://www.milenio.com/policia/jose-francisco-ruiz-massieu-quien-lo-mato-y-por-que}{Francisco Ruiz Massieu}, and drug related coverage such as the assassination of \href{https://www.aciprensa.com/Docum/posadas.htm}{Cardinal Juan Jesus Ocampo}.}

The efforts of independent media, combined with the establishment of the autonomous electoral monitoring organization, the Federal Electoral Institute (\textit{IFE}, by its Spanish acronym), paved the way for significant political changes in Mexico. These included left-wing victories in several states, most notably Cuauhtémoc Cárdenas's win in Mexico City in 1997 and Andrés Manuel López Obrador's (AMLO) victory in 2000. Additionally, this change in behavior by the media contributed to the end of the PRI's 70-year hold on the presidency with Vicente Fox's historic win for the National Action Party (PAN, by its acronym in Spanish) \citep{Lawson-AMLO}.

\subsubsection*{PAN and the War on Drugs (2000-2012)}

Although the democratic transition in Mexico is thought to have improved the country's democratic functioning, some other significant problems arised. The most prominent one, drug cartel organizations (DTOs) and drug cartel related violence. During the period of 2006 to 2020, Mexico experienced alarmingly high and escalating levels of violence due in part to the so called \textit{War on Drugs}, where the Mexican federal government launched a massive military and police campaign against DTOs. This intervention by President Felipe Calderon aimed to end twelve years of inter-cartel wars over control of drug trafficking routes into the United States. The campaign notably back-fired, leading to a 174\% inter-annual increase on intentional homicides per 100 thousand inhabitants and to the expansion and re-organization of DTOs into new illegal markets—such as extortion, kidnapping for ransom, human smuggling, among others. \footnote{According to \citep{Patrones-INEGI}, the rate of homicides per 100 thousand inhabitants increased from 8.1 on 2007 to 22.2 on 2012. Impressively, from 2006 to 2012 the number of missing persons increased a staggering 1,473\% \href{https://es.statista.com/estadisticas/1268415/numero-anual-de-personas-desaparecidas-en-mexico/}{Statista}} 

Drug-related violence was an issue of significant electoral importance for the PAN, particularly because the party's decision to initiate a war on drugs led to an exponential increase in violence. To mitigate potential electoral consequences—a strategy that ultimately failed, as the PAN lost the 2012 election—they resorted to some of the old practices of manipulating the information environment. Evidence of this issue perceived salience, is that this surge in violence was not uniform across the country; it was especially severe in regions governed by opposition parties. In these areas, the federal government allocated fewer resources to combat violence, exacerbating the problem \citep{Federalism-Trejo-PoliGob}.

The Florence Cassez case is a notorious example of government-media collusion in Mexico related to violence related issues. In 2005, French national Florence Cassez and her boyfriend, Israel Vallarta, were accused of running a kidnapping ring. Their \textit{live} arrest and the dramatic rescue of victims were broadcast by major networks Televisa and TV Azteca, but it was later revealed that the operation had been staged by law enforcement a day after their actual arrest \citep{Cassez-Guardian}. This revelation exposed the manipulation of public perception through media, raising concerns about the integrity of Mexico's justice system and its ties to the press. Cassez was convicted and sentenced to 60 years in prison, but in 2013, the Mexican Supreme Court released her, citing serious due process violations \citep{Cassez-BBC}. The case remains a symbol of corruption, media sensationalism, and the flaws in Mexico's legal system. 

Calderón's relationship with the media was initially presented as one of cooperation, as he suggested that media coverage of violent events could inadvertently benefit drug trafficking organizations (DTOs) by revealing government plans and strategies. On 2011, one year before the elections, Calderón met with representatives from Mexico's largest media outlets, convened by Televisa and TV Azteca, to establish an agreement on how drug violence would be reported \citep{Callaron-Jornada}. In this meeting, the government agreed to provide protection for journalists in exchange for their adherence to editorial guidelines that would '\textit{not benefit criminal organizations}' or portray them as positive contributors to society. Years later, it was revealed that 715 media outlets signed Calderón's proposal to be part of his strategy and '\textit{self-regulate}' their content. One ambiguous condition in the agreement stated: '\textit{Not to interfere in the fight against crime}' \citep{Callaron-Proceso}.

Later, it also came to light that the government was illegally channeling funds from the Ministry of Health for advertising purposes. In one instance, approximately 1.5 billion Mexican pesos were diverted from health programs for this purpose \citep{Contralinea-Desvio}. Moreover, during 2009, spending on advertising under Calderón's administration increased by 501\% compared to the last year of Vicente Fox's presidency. The amount spent in 2009 exceeded the budget for that year by 145\% and was 49\% higher than the amount spent in 2008. These financial practices underscored the administration's heavy reliance on media spending to control narratives and shape public perception \citep{Fundar}.

\subsubsection*{Return of PRI (2012-2018)}

On 2012, Enrique Peña Nieto from PRI party won the elections against Josefina Vasquez Mota (PAN) and Andres Manuel Lopez Obrador (PRD). While on campaign, President Enrique Peña Nieto promised to steer away from his predecessor's policies, however, his administration saw increased militarization, marked by blatant incompetence and human rights violations.\footnote{The most relevant ones being the enforced disappearance of 43 Ayotzinapa students, in collusion with the federal government and the army, as well as the extrajudicial executions, including those in Tlatlaya, Apatzingán and Tanhuato \citep{Amnesty-Int}} Consequently, he concluded his tenure with the highest ever rate of intentional homicides, reaching 30 per 100,000 inhabitants annually on 2018 \citep{Intentional-Homicides}.

Evidence of the strategic use of advertising spending during this period is well-documented. For instance, in 2017, The New York Times published an article titled "\textit{Using Billions in Government Cash, Mexico Controls News Media}" \citep{Times}. The article highlights how running a media outlet in Mexico often depended on a single client—the government—which paid for advertising only on the condition that the outlet refrained from criticizing it. The report is striking, as it provides accounts from journalists, editors, and managers who revealed that government officials routinely dictated coverage, instructing outlets on what they should—and should not—report. Investigative stories were frequently softened, delayed, or entirely suppressed if they were covered at all. For example, one of the most high-profile cases involved the removal of a Milenio story that criticized the government’s anti-poverty program for misusing funds after a meeting with the Secretary of Social Development, who served three years in prison for corruption scandals \citep{Milenio-Buzz}.

Other tactics that I documented from PRI before were also employed. Undisclosed firings of journalists were common, for example the firing from MVS of news reporter Carmen Aristegui after she published the infamous corruption scandal of \textit{The White House} \citep{WashPost-Aristegui}. In another article, \cite{Times2-Kill} presents further evidence that PRI tactics of intimidation and violence, including murder, persisted during Peña Nieto's administration. One chilling example cited is the death of a Proceso photographer who was killed in his Mexico City apartment after fleeing Tabasco, where he had been harassed by local authorities. These examples underline how both financial control and violent repression were once again used to maintain influence over the media and suppress dissenting voices.

\subsubsection*{AMLO and MORENA (2019-2024)}

In 2018, Andrés Manuel López Obrador (AMLO), the candidate for the newly formed MORENA party, won the presidential election, signaling a significant shift in government-media relations. Upon taking office, AMLO inherited a nation grappling with escalating violence. Despite campaign promises to reduce militarization and address the root causes of insecurity, his early years in office were marked by continued high levels of violence, with intentional homicides and other crimes reaching unprecedented levels. While the creation of the National Guard aimed to address these issues, the outcomes were widely criticized, as the violence remained persistent despite stopping its inter-annual increases.

Under his administration’s policy of republican austerity, government advertising spending was reduced by 78.7\%, likely influencing media coverage and their stance on various issues \citep{Article19, ElEconomist-Republican}. Figure \ref{fig:yearly_spending} illustrates this sharp decrease in spending. However, this reduction did not lead to a more transparent or equitable allocation of resources. \cite{Article19} notes that the absence of a robust regulatory framework allowed the distribution of funds to remain highly discretionary. Investigative reports revealed a shift in spending away from right-wing opposition outlets toward media outlets aligned with the administration. For example, Nexos, a right-wing opinion magazine, lost its entire advertising budget after AMLO took office, while southern outlets like \textit{Grupo Cantón} and \textit{Compañía Editorial del Mayab}, which maintained close ties with AMLO, saw significant increases in funding \citep{Animal-Discrecionalidad}. Nonetheless, the nearly 80\% reduction meant that in general all large outlets experienced cuts.

The question of why these cuts were introduced, despite their potential electoral costs, is undoubtedly puzzling. One possible explanation is that the move was intended to gain favor with the electorate by exposing the substantial sums of money media outlets received under previous administrations, thereby delegitimizing these outlets. This aligns with his broader narrative of promising to reduce spending on such expenditures to reallocate funds to social programs and other critical government priorities.

Another potential explanation, which does not entirely conflict with the first, is more strategic: cutting the significant payments that media outlets had long been used to may have pushed them to the brink of financial instability. This is actually really credible as \cite{Times} documents, the government was the biggest client of most of these outlets. This approach would have allowed the administration to exert greater control over them throughout his tenure, adjusting advertising spending marginally as a way to manage media coverage, particularly during moments of political vulnerability, such as major corruption scandals.

In addition, López Obrador introduced daily morning press conferences, where he engaged directly with the media every weekday. These sessions were used to address pressing issues, express his thoughts, criticize opposition parties and figures, and respond to questions from the press and public. This unprecedented level of accessibility reshaped the interaction between the government and the media, creating a new dynamic in political communication \citep{ThugMex-VictorHugo-JourPract}.

While these press conferences could have served as an accountability mechanism, many critics argue they functioned primarily as platforms for government propaganda. Scholars and journalists contend that AMLO used them to spread misleading narratives, consolidate support for his administration, and publicly attack dissenters and opposition figures \citep{OtrosDatos}.


